%=============================================================================!
% 18.7  MOMENTUM AND ADAM                                                     !
%=============================================================================!
\section{Momentum and Adam}
\label{sec:lo-momentum}

A heavy block sliding down a long, gently tilted gutter with steep sides
gathers speed along the gutter, since every push points the same way
there, while across it the pushes alternate from one wall to the other
and friction damps the swaying. The walker of \cref{sec:lo-descent} has
no memory: each step depends only on the slope underfoot, so along a
long shallow valley it creeps, and across a steep one it zig-zags. Giving
the steps memory, as the block's speed is a memory of past pushes, cures
much of both.

\subsection*{Momentum}

\emph{Gradient descent with momentum} keeps a velocity $\vb{v}$, a
running sum of past gradients that decays by a factor $\mu < 1$ a step,
and moves the weights along it:
\begin{equation}
   \vb{v} \leftarrow \mu\vb{v} + \grad L(\vb{w}) , \qquad
   \vb{w} \leftarrow \vb{w} - \eta\vb{v} .
   \label{eq:lo-momentum}
\end{equation}
The block is more than a picture. Writing $\vb{w}_t$ for the weights after
$t$ steps, the second rule gives $\eta\vb{v}_t = \vb{w}_{t-1} - \vb{w}_t$
at every step. Multiplying the first rule by $-\eta$ and substituting
gives
\begin{equation*}
   \vb{w}_{t+1} - \vb{w}_t = \mu(\vb{w}_t - \vb{w}_{t-1})
   - \eta\grad L(\vb{w}_t) ,
\end{equation*}
and moving $\vb{w}_t - \vb{w}_{t-1}$ to the left,
\begin{equation*}
   (\vb{w}_{t+1} - 2\vb{w}_t + \vb{w}_{t-1})
   + (1 - \mu)(\vb{w}_t - \vb{w}_{t-1}) = -\eta\grad L(\vb{w}_t) .
\end{equation*}
The first bracket is a second difference, the change of the step from
one step to the next, an acceleration, and the second a velocity: this is
Newton's law with friction, $m\ddot{\vb{w}} + \gamma\dot{\vb{w}} =
-\grad L$ (\cref{sec:os-damped}), discretised with a step of one, the
mass one, the friction $1 - \mu$ and the force $\eta$ times the gradient.
Momentum turns the descent into the damped motion of a particle on the
loss.

On a quadratic loss each eigen-direction of the Hessian again moves on
its own, the component $c_t$ along an eigenvector of eigenvalue
$\lambda$ following $c_{t+1} = (1 + \mu - \eta\lambda)\,c_t - \mu
c_{t-1}$, the line above with $\grad L$ replaced by $\lambda c_t$.
Trying $c_t = z^t$, as for the oscillations of \cref{sec:in-stability},
\begin{equation*}
   z^2 - (1 + \mu - \eta\lambda)\,z + \mu = 0 .
\end{equation*}
The product of the two roots is $\mu$, so when they are complex, a pair
of conjugates of equal size, both have $\lvert z\rvert = \sqrt\mu$: the
component swings and shrinks by $\sqrt\mu$ a step, whatever $\lambda$.
They are complex while $(1 + \mu - \eta\lambda)^2 \le 4\mu$, and the
best choice, due to Polyak\cite{Polyak1964}, puts both ends of the range
of eigenvalues on the edge of this condition, $1 + \mu -
\eta\lambda_{\min} = 2\sqrt\mu$ and $1 + \mu - \eta\lambda_{\max} =
-2\sqrt\mu$:
\begin{align*}
   (1 - \sqrt\mu)^2 &= \eta\lambda_{\min} , \quad
   (1 + \sqrt\mu)^2 = \eta\lambda_{\max}
      && \text{as squares}\\
   \frac{1 + \sqrt\mu}{1 - \sqrt\mu} &= \sqrt\kappa
      && \text{dividing, and the root}\\
   \sqrt\mu &= \frac{\sqrt\kappa - 1}{\sqrt\kappa + 1}
      && \text{solving for $\sqrt\mu$,}
\end{align*}
the first line writing $1 + \mu \mp 2\sqrt\mu$ as $(1 \mp \sqrt\mu)^2$.
Adding the square roots of the first line's two equations gives $2 =
\sqrt\eta\,(\sqrt{\lambda_{\min}} + \sqrt{\lambda_{\max}})$, so $\eta =
4/(\sqrt{\lambda_{\max}} + \sqrt{\lambda_{\min}})^2$.
The asymptotic contraction factor in every direction is $(\sqrt\kappa - 1)/(\sqrt\kappa + 1)$ a step,
against $(\kappa - 1)/(\kappa + 1)$ without momentum, and the steps
needed grow as $\sqrt\kappa$ instead of $\kappa$. For the valley of
\cref{fig:lo-optimisers}, with eigenvalues 1 and 100, the factors are
0.818 and 0.980: momentum brings the loss below $10^{-8}$ in 77 steps,
while gradient descent at its best fixed rate has not in 300. The steep
weight shows the difference (\cref{fig:lo-optimisers}a): without
momentum it flips sign at every step and barely shrinks, by the factor
$-0.98$; with momentum it rings and dies away, the damped swing of
Chapter~4.

\subsection*{Adam}

The networks of the later chapters are trained with
Adam\cite{Kingma2014}, which gives each weight a step of its own size.
It keeps two running means of each weight's gradient, of the gradient
and of its square,
\begin{equation*}
   m \leftarrow \beta_1m + (1 - \beta_1)\,g , \qquad
   s \leftarrow \beta_2s + (1 - \beta_2)\,g^2 ,
\end{equation*}
with $\beta_1 = 0.9$ and $\beta_2 = 0.999$ by default, each a mean over
roughly the last $1/(1 - \beta)$ steps, 10 and 1000, weighted towards the
recent ones. Both start at zero, which biases them low at first: with a
steady gradient $g$, after $t$ steps
\begin{equation*}
   m_t = (1 - \beta_1)\sum_{k=0}^{t-1}\beta_1^kg
   = (1 - \beta_1)\frac{1 - \beta_1^t}{1 - \beta_1}\,g
   = (1 - \beta_1^t)\,g ,
\end{equation*}
by the sum of a geometric series, so dividing by $1 - \beta_1^t$ undoes
the bias, and likewise for $s$. The step is
\begin{equation*}
   w \leftarrow w - \eta\,\frac{\hat m}{\sqrt{\hat s} + \epsilon} , \qquad
   \hat m = \frac{m}{1 - \beta_1^t} , \quad \hat s = \frac{s}{1 - \beta_2^t} ,
\end{equation*}
with $\epsilon = 10^{-8}$ guarding against a division by zero. For a
steady gradient $\hat m/\sqrt{\hat s}$ is $\pm1$, so each weight moves by
about $\eta$ a step whatever the size of its gradient; multiplying one
weight's gradients by a constant multiplies $\hat m$ and $\sqrt{\hat s}$
alike and leaves its steps unchanged. In \texttt{mdlab}:

\begin{code}
class Adam:
    def __init__(self, rate: float = 1e-3, beta1: float = 0.9,
                 beta2: float = 0.999, eps: float = 1e-8):
        self.rate, self.beta1, self.beta2, self.eps = rate, beta1, beta2, eps
        self.m = self.s = None
        self.t = 0

    def step(self, w: NDArray, g: NDArray) -> NDArray:
        if self.m is None:
            self.m, self.s = np.zeros_like(g), np.zeros_like(g)
        self.t += 1
        self.m = self.beta1 * self.m + (1 - self.beta1) * g
        self.s = self.beta2 * self.s + (1 - self.beta2) * g * g
        m_hat = self.m / (1 - self.beta1**self.t)
        s_hat = self.s / (1 - self.beta2**self.t)
        return w - self.rate * m_hat / (np.sqrt(s_hat) + self.eps)
\end{code}

\noindent A test checks that its steps, and those of \texttt{Momentum},
match PyTorch's \texttt{Adam} and \texttt{SGD} with momentum on a loss of
two weights to the last digit.

The scaling is Adam's strength and its limit. A valley whose steep and
shallow directions lie along the weights, as when one weight multiplies
a feature a hundred times larger than another's, is a valley of badly
scaled weights, and Adam evens them out: with the rate 0.05 it reaches
$10^{-8}$ in 189 steps (\cref{fig:lo-optimisers}b). The same valley
turned by \ang{45}, its steep direction a mixture of both weights, gives
each weight a similar gradient, Adam's scaling has nothing to even out,
and after 300 steps its loss is still 0.49, worse than plain descent's
\num{0.0014}, while momentum still needs only 85
(\cref{fig:lo-optimisers}c). Adam is robust where the scales of the
weights differ widely and are not known in advance, as in a network,
which is why it is the common default; a narrow valley at an angle it
leaves as it found it.

\begin{figure}[htb]
\centering
\includegraphics{ch18_learning/optimisers}
\caption{Gradient descent, momentum and Adam in a valley whose Hessian
has the eigenvalues 1 and 100, from $(-2.5, 0.5)$: descent at its best
fixed rate $2/(\lambda_{\min} + \lambda_{\max})$, momentum with Polyak's
rate and $\mu$, Adam with the rate 0.05. (a) The steep weight against
the step, the valley's axes along the weights. (b) The loss against the
step there, with the factors per step that theory gives descent and
momentum (dashed). (c) The loss in the same valley turned by \ang{45}.}
\label{fig:lo-optimisers}
\end{figure}

\begin{keyidea}
Momentum adds to the step a decaying memory of past gradients; it is
damped motion on the loss, and on a quadratic with Polyak's choice it
has the asymptotic contraction factor $(\sqrt\kappa - 1)/(\sqrt\kappa + 1)$ a step.
Adam divides each weight's running mean gradient by the root of its
running mean square, with both corrected for their start at zero, so that
every weight moves at about the rate whatever the scale of its gradient;
it evens out badly scaled weights, and leaves a narrow valley at an angle
as narrow as it was.
\end{keyidea}

\notebook{18}{race the three optimisers down a valley, and turn the valley}

\begin{selfcheck}
For a Hessian with $\kappa = 10^4$, by what factor does the slow direction
shrink in a step of gradient descent at its best fixed rate, and of
momentum with Polyak's choice?
\end{selfcheck}
